\begin{longtable}{@{}P{1.55cm}P{2.9cm}P{4.5cm}P{3.2cm}P{3.4cm}P{4.6cm}P{3.6cm}@{}}
\caption{Data extraction (part 1). `Not reported' marks items the paper does not report.}\label{tab:ext1}\\
\toprule Study & Problem / aim & Method & Dataset(s) & Evaluation & Key findings & Venue; peer review \\ \midrule \endfirsthead
\toprule Study & Problem / aim & Method & Dataset(s) & Evaluation & Key findings & Venue; peer review \\ \midrule \endhead
\bottomrule \endfoot
\textbf{S01} \cite{S01} Alvarez-Melis \& Jaakkola (2018); RQ1, RQ2 & Robustness (similar inputs should get similar explanations) is a key desideratum of interpretability but is not measured. & Local Lipschitz estimate of explanation change within an eps-ball (continuous optimisation; discrete sample-based version for categorical data); LIME, SHAP and gradient-based methods. & UCI datasets (e.g., Boston) and COMPAS with black-box classifiers; MNIST and ImageNet for gradient methods (p.3). & Local Lipschitz estimates aggregated over test points. & Current methods, including LIME and SHAP, performed poorly on the robustness metrics; very close points received considerably different explanations (abstract; Sec.3). & ICML 2018 Workshop on Human Interpretability in Machine Learning (WHI); arXiv:1806.08049; Workshop. Workshop paper (arXiv comment) \\
\textbf{S02} \cite{S02} Zhang et al. (2019); RQ1, RQ3 & LIME explanations contain uncertainty that can undermine users' trust in predictions. & Repeated LIME runs (100 per point) on a random forest trained on synthetic tree-generated data; varying sample size and sampling proximity; text and COMPAS case studies. & Synthetic 8- and 4-feature tree-generated data; 20 Newsgroups; COMPAS (pp.2-6). & Empirical selection probability of features across repeated LIME runs. & Sampling randomness, the sample-size and proximity settings, and varying model credibility across points all create uncertainty; LIME did not always select the locally important features (Fig.2, p.3). & ICML 2019 Workshop on AI for Social Good; arXiv:1904.12991; Workshop. Workshop paper (PDF header) \\
\textbf{S03} \cite{S03} Saito et al. (2020); RQ3 & LIME's Gaussian sampling creates off-manifold samples that an adversary can exploit to hide biased behaviour. & CTGAN-based locality sampling inside LIME (CTGAN-LIME) vs default LIME; random-forest black box; top-k detection of the hidden sensitive feature. & Three real-world tabular datasets including COMPAS and Communities and Crime (pp.1-4). & Accuracy of detecting the biased feature among the top-k explanation features. & CTGAN-LIME detected the hidden biased behaviour more accurately than LIME across three datasets, up to 99.94\% top-1 accuracy in some cases (abstract). & AdvML '20: Workshop on Adversarial Learning Methods for Machine Learning and Data Mining...; Workshop. Workshop paper (PDF header) \\
\textbf{S04} \cite{S04} Slack et al. (2020); RQ3 & Perturbation-based explainers can be fooled because their synthetic samples lie off the data manifold. & Scaffolded classifier with an out-of-distribution detector: biased model on real data, innocuous model on LIME/\allowbreak{}SHAP perturbations. & COMPAS, Communities and Crime, German Credit (p.2). & Share of explanations in which the sensitive feature appears at each rank; OOD-detector fidelity. & LIME and SHAP explanations of the biased classifiers did not flag the sensitive attribute; LIME was more vulnerable than SHAP (abstract; Conclusion). & Proceedings of the AAAI/\allowbreak{}ACM Conference on AI, Ethics, and Society; Conference. Peer-reviewed conference proceedings (Crossref) \\
\textbf{S05} \cite{S05} Slack et al. (2021); RQ1, RQ3 & LIME and KernelSHAP explanations are inconsistent and unstable, give no measure of their reliability and need hyper-parameter tuning. & Bayesian versions of LIME and KernelSHAP that return credible intervals for feature importances; 'perturbations-to-go' estimate of the samples needed for a target certainty; focused sampling of perturbations. & COMPAS, German Credit, ImageNet and MNIST (p.2). & Calibration of credible intervals; stability by local Lipschitz estimate (Eq.12) with Wilcoxon signed-rank tests (pp.8-9). & BayesLIME and BayesSHAP were significantly more stable than LIME and SHAP (about 53\% on average, p \textless{} 0.01), except BayesSHAP on German Credit (p.9). & Advances in Neural Information Processing Systems 34 (NeurIPS 2021); Conference. Peer-reviewed conference (Semantic Scholar) \\
\textbf{S06} \cite{S06} Yasodhara et al. (2021); RQ1 & Accuracy and stability of global feature-importance methods (gain, TreeSHAP) for tree ensembles had not been compared. & Random forest, gradient boosting (sklearn) and XGBoost explained with gain and path-dependent TreeSHAP (global importances). Input perturbation: Gaussian noise with SD of 0.5x, 1x or 2x each feature's SD. Model perturbation: retraining with different seeds/\allowbreak{}hyperparameters while keeping predictions highly correlated. 50 repetitions (5 for one dataset). & Synthetic classification (300 samples, 5-150 features); Forest Fire (517 x 12), Concrete (1,030 x 8), Auto MPG (406 x 7), proprietary Company Finance (2,716 x 892) (Table 1, p.5). & Accuracy of top-feature ranking against known synthetic coefficients; stability of importance rankings across perturbed inputs/\allowbreak{}models; comparison gain vs SHAP across three ensembles. & Global importance methods lack accuracy with and without noisy inputs and become unstable as input dimension or noise increases and when models are re-initialised with different seeds or hyperparameters (abstract, Sec.4). & Machine Learning and Knowledge Extraction (CD-MAKE 2021), Lecture Notes in Computer...; Conference. Peer-reviewed conference proceedings (Crossref) \\
\textbf{S07} \cite{S07} Zafar \& Khan (2021); RQ1, RQ3 & LIME's random perturbation gives different explanations for the same prediction. & DLIME: agglomerative hierarchical clustering and KNN select a cluster; a linear model or tree trained on it gives deterministic explanations; stability measured as variability of selected feature sets (Jaccard) over repeated runs. & Three binary and three multi-class public datasets plus six synthetic datasets. & Jaccard-based stability over repeated explanations; faithfulness. & DLIME was deterministic and more stable than LIME while remaining faithful (abstract). & Machine Learning and Knowledge Extraction; Journal. Peer-reviewed journal (Crossref) \\
\textbf{S08} \cite{S08} Zhou et al. (2021); RQ1, RQ3 & Random generation of perturbation samples makes LIME explanations unstable across runs. & S-LIME: hypothesis tests based on a central limit theorem for LARS/\allowbreak{}LASSO variable-entry decisions decide adaptively how many perturbation samples LIME needs for a stable top-k feature set. & Simulated data; Breast Cancer Wisconsin (569 x 30; random forest, 500 trees, \textasciitilde{}95\% test accuracy); a sepsis-prediction RNN on temporal EHR data (p.5). & Average pairwise Jaccard index of top-k features (k=1..5) over 20 repetitions; baseline: standard LIME with 1,000 perturbations; nmax = 10,000 for S-LIME (Table 1, p.5). & On a breast-cancer test instance, LIME's average Jaccard was 0.61 (k=1), 0.66 (k=4) and 0.59 (k=5), versus 1.0, 1.0 and 0.85 for S-LIME (Table 1). & Proceedings of the 27th ACM SIGKDD Conference on Knowledge Discovery \& Data Mining; Conference. Peer-reviewed conference proceedings (Crossref) \\
\textbf{S09} \cite{S09} Agarwal et al. (2022); RQ1, RQ2 & Existing stability metrics consider only input change and ignore the model's internal representation and output. & Proposes Relative Input, Representation and Output Stability (RIS, RRS, ROS) computed over label-preserving perturbations (\textbar{}\textbar{}x-x'\textbar{}\textbar{} \textless{}= eps, same predicted class); proves RIS \textless{}= lambda * L1 * RRS. & Adult, COMPAS, German Credit (tabular); ANN and logistic-regression predictors (Fig.3, App.). & RIS/\allowbreak{}RRS/\allowbreak{}ROS (log scale) for Gradient, IntegratedGrad, Input*Grad, SmoothGrad, SHAP, LIME and a random baseline; empirical vs theoretical bounds. & No method is completely stable; SmoothGrad is most stable on RRS/\allowbreak{}ROS, beating others by 12.7\% on average; theoretical bounds are on average 233\% above empirical values (p.5). & ICLR 2022 Workshop (non-archival); arXiv:2203.06877; Workshop. Workshop paper (non-archival); arXiv version \\
\textbf{S10} \cite{S10} Visani et al. (2022); RQ1, RQ2, RQ3 & LIME gives different explanations for the same instance across repeated calls; the paper proposes indices to detect this instability. & Derives the sampling distribution of LIME's weighted ridge coefficients and proposes two indices: Variables Stability Index (concordance of selected features) and Coefficients Stability Index (equality tests on coefficients) across repeated LIME calls. & Credit-scoring data from CRIF S.p.A. (about 100 variables, top 20 retained; sample size not reported) (Sec.6, p.12). & VSI and CSI (0-100) over 10 repeated LIME calls using the 7 most important features; logistic regression benchmark vs gradient boosting trees; kernel width and ridge penalty varied (pp.14-15). No external baseline stability metric. & Stability depends strongly on LIME hyperparameters: for the same applicant VSI 89.44\% /\allowbreak{} CSI 92.7\% (kernel width 3, ridge 1) versus VSI 14.17\% /\allowbreak{} CSI 57.46\% (kernel width 1.3, ridge 0.001) (p.15). High dimensionality erodes LIME locality. & Journal of the Operational Research Society; Journal. Peer-reviewed journal (Crossref) \\
\textbf{S11} \cite{S11} Ahn et al. (2023); RQ3 & When only a static sample of model inputs and outputs is available, local explanations carry finite-sample uncertainty that current explainers do not quantify. & Model-agnostic local explanations built from a fixed input-output sample (local polynomial regression, plug-and-play with other local explainers); non-parametric bootstrap uncertainty intervals for feature-importance scores; compared with naive regression confidence intervals and BayesLIME. & Simulated ground-truth function; COMPAS and German Credit with random forests; Spambase with a SAINT neural network (p.8). & Average interval width and empirical coverage probability in simulation; visual summaries on real data. & Bootstrap intervals gave a better width-coverage trade-off than naive intervals and higher coverage than BayesLIME without model access (pp.1,8). & arXiv:2301.05761; Preprint. Preprint (not peer reviewed) \\
\textbf{S12} \cite{S12} Dhurandhar et al. (2023); RQ1 & LIME and variants give explanations that change across nearby examples and sampling runs. & LINEX: local explanations from an invariant-risk-minimisation game over several perturbation environments, eliminating features whose gradient flips sign locally; random and realistic (MAPLE-like) neighbourhoods. & Tabular (e.g., IRIS, MEPS and others in Table 3) and image (FMNIST, CIFAR10) datasets; 80/\allowbreak{}20 splits. & Infidelity, generalised infidelity, Coefficient Inconsistency (CI), Class Attribution Consistency (CAC), unidirectionality; baselines LIME, S-LIME, MeLIME, NAS-LIME; paired t-tests (p.9). & LINEX was better than baselines in 21 of 40 cases and worse in 5, with notably better stability metrics (CI, CAC) and unidirectionality (p.9). & Advances in Neural Information Processing Systems 36; Conference. Peer-reviewed conference proceedings (Crossref) \\
\textbf{S13} \cite{S13} Pawlicki (2023); RQ1, RQ2 & Measures of explanation quality are lacking; how stable are SHAP explanations under data perturbation? & Random forest; three perturbation types (feature shuffling, missing values, Gaussian noise) at four strengths; SHAP values from models trained on clean data and applied to perturbed data; correlation and error measures; Shapiro-Wilk then paired t-test or Wilcoxon. & Breast Cancer, Iris, Wine and CIC-IDS2018 (p.2). & Correlation and error between clean and perturbed SHAP values; statistical tests of distribution equality (p.5). & SHAP values changed with perturbation strength; stability varied by dataset, with CIC-IDS2018 showing the largest decreases (pp.10). & 2023 IEEE 10th International Conference on Data Science and Advanced Analytics (DSAA); Conference. Peer-reviewed conference proceedings (Crossref); Zenodo pre-publication version read \\
\textbf{S14} \cite{S14} Yuan et al. (2023); RQ1, RQ3 & Effect of background-dataset size on the stability of SHAP values and rankings for neural-network models. & Three-layer ANN on mortality prediction; background sets of 100-1,000 instances sampled 100 times each; SHAP values and model-level rankings recomputed on the explanation set. & MIMIC-III: 44,918 ICU admissions (3,958 positive); 70/\allowbreak{}30 development/\allowbreak{}explanation split (p.4). & Variance of instance-level SHAP values; quartile-based BLEU and Jaccard (proposed) for ranking stability; no alternative explainer baseline. & SHAP values and rankings fluctuate under small random background sets; fluctuation decreases as background size grows; a U-shape shows top- and bottom-ranked variables are more stable than mid-ranked ones (abstract, Sec.3). & ICLR 2023 Tiny Papers track; arXiv:2204.11351; Workshop. Peer-reviewed workshop track (per DBLP and Semantic Scholar) \\
\textbf{S15} \cite{S15} Badhon et al. (2024); RQ1, RQ3 & LIME's random neighbourhood sampling causes selection variability and class imbalance in neighbourhoods, reducing stability. & CTGAN-LIME: CTGAN-generated neighbourhood samples with conditional vectors for class balance, plus real samples; ridge surrogate; feedforward neural-network black box. & Heart Disease, Red Wine Quality and Census Income (Fig.3, p.7). & Local fidelity, stability and local concordance vs LIME. & CTGAN-LIME improved reliability and stability of explanations relative to LIME across the datasets (abstract; Fig.3). & World Congress on Electrical Engineering and Computer Systems and Science; Conference. Peer-reviewed conference proceedings (Crossref) \\
\textbf{S16} \cite{S16} Baniecki \& Biecek (2024); RQ3 & Explanations can be manipulated, fooled or fairwashed; a unified taxonomy of attacks and defences is lacking. & Narrative survey with unified notation and taxonomy of ten attack types and defences; tabulates prior works by dataset, model and explanation type. & Not applicable (review of published studies; tabulated studies include tabular datasets such as Adult and Heart). & Not applicable (review). & Post-hoc explanations (incl. SHAP/\allowbreak{}LIME, PDP) are vulnerable to input, model and data manipulation; robust interpretation and evaluation protocols must account for adversarial settings (abstract; Sec.6). & Information Fusion; Journal. Peer-reviewed journal (Crossref) \\
\textbf{S17} \cite{S17} Elgezawy et al. (2024); RQ1, RQ3 & Randomised synthetic sampling makes LIME explanations unstable. & ST-LIME: fixed-noise synthetic sample generation scaled by feature standard deviations; compared with five state-of-the-art LIME variants over 10 repeated explanations per sample. & Parkinson's, Breast Cancer, EEG and Bank datasets (p.7). & Feature Sequence Stability Index (FSSI) and fidelity. & ST-LIME reached FSSI = 1.0 on the datasets reported and outperformed the compared techniques on average (Table 2, p.8). & International Journal of Intelligent Engineering and Systems; Journal. Peer-reviewed journal (Crossref) \\
\textbf{S18} \cite{S18} Goldwasser \& Hooker (2024); RQ3 & Sampling approximations of Shapley values carry considerable Monte-Carlo uncertainty. & ControlSHAP: control variates based on Taylor (linear/\allowbreak{}quadratic) approximations of the model for both independent and correlated feature settings. & Simulated correlated data; UCI Portuguese bank, German credit, census income and breast-cancer subtyping (p.9). & Variance reduction of Shapley estimates and number of ranking changes across repetitions. & Final Shapley estimates had roughly 80-95\% lower variability than the original estimates; after 100 iterations ControlSHAP matched the variability of 1,000 plain iterations (pp.10,14). & Explainable Artificial Intelligence (xAI 2024), Part II, Communications in Computer and...; Conference. Peer-reviewed conference proceedings (Crossref) \\
\textbf{S19} \cite{S19} Hill et al. (2024); RQ2, RQ3 & Local explanations near complex decision boundaries are brittle and their uncertainty is not quantified. & Gaussian Process Explanation UnCertainty (GPEC): GP regression over explanations with a decision-boundary-aware geodesic kernel, combining boundary complexity with function-approximation uncertainty; KernelSHAP by default. & Three tabular datasets (Census, Online Shoppers, German Credit), synthetic data and image datasets (p.6). & Comparison with BayesSHAP, BayesLIME, CXPlain and a naive GP; capture of boundary complexity; ablation; hyperparameter and sample-size sensitivity (appendix). & GPEC uncertainty grew with decision-boundary complexity and gave better insight into explanation reliability than existing methods (abstract; Sec.5). & Proc. 27th International Conference on Artificial Intelligence and Statistics (AISTATS...; Conference. Peer-reviewed conference (PDF footer; arXiv comment) \\
\textbf{S20} \cite{S20} Ihalapathirana et al. (2024); RQ1, RQ2, RQ3 & Effect of data size and class imbalance on model performance and SHAP stability. & Linear discriminant analysis, XGBoost and a neural network on five predictors; varied training/\allowbreak{}background data sizes and balance; 10-fold cross-validation; UMAP of mean \textbar{}SHAP\textbar{} across folds. & North Karelia Wellbeing Services County electronic health records for type-2 diabetes (7,601 samples after de-duplication; 500 held-out test) (p.2). & Variance of SHAP values across folds/\allowbreak{}background sets; model performance metrics. & Larger and balanced background data reduced SHAP variance; data characteristics that maximise performance did not always give stable SHAP explanations (abstract; p.5). & ESANN 2024 proceedings; Conference. Peer-reviewed conference proceedings (Crossref) \\
\textbf{S21} \cite{S21} Kelodjou et al. (2024); RQ1, RQ3 & Kernel SHAP gives different explanations for the same input across executions because of stochastic neighbour (coalition) selection. & ST-SHAP: deterministic, layer-wise neighbour selection (layer-1 coalitions) that removes sampling instability while keeping fidelity; theoretical analysis of layer-1 attribution scores. & Tabular datasets incl. Adult, Dry Bean, Movie and Credit Card (Table 4). & Jaccard coefficient of non-zero feature sets across repeated executions; Kendall's tau and R2 vs exact SHAP values. & ST-SHAP achieved full stability without compromising fidelity; layer-1 scores stayed close to exact SHAP values (Table 4; Conclusion). & Proceedings of the AAAI Conference on Artificial Intelligence; Conference. Peer-reviewed conference proceedings (Crossref) \\
\textbf{S22} \cite{S22} Ketata et al. (2024); RQ1, RQ2 & Explanations can vary considerably for similar instances and across data samples, undermining reliability in healthcare. & SHAP and LIME computed within k-fold cross-validation for several k; random forest; Spearman similarity of feature-rank lists across k (generalisability), concordance with internal RF importance, and a stability metric. & Hypothyroidism and diabetes risk-prediction datasets (p.7). & Spearman correlation of rankings across k values; concordance; stability. & SHAP with k-fold showed higher generalisability, stability and concordance than LIME; LIME showed poor concordance and reliability (pp.13-14). & IEEE Access; Journal. Peer-reviewed journal (Crossref) \\
\textbf{S23} \cite{S23} Khan et al. (2024); RQ2, RQ3 & Explainers should give similar explanations for similar inputs, yet small input changes can shift explanations. & Defines explainer astuteness and proves lower bounds linking it to the predictor's probabilistic Lipschitzness for SHAP, RISE and CXPlain; four classifiers per dataset; 20 explainer runs. & Simulated Orange Skin, Nonlinear Additive and Switch; UCI Rice and Telescope; MNIST (p.7). & Empirical astuteness against the classifiers' estimated Lipschitz constants; comparison across explainers. & Smoother (lower-Lipschitz) predictors gave more astute explanations; RISE was least and CXPlain most astute on the evaluated data (pp.8-9). & Proc. 27th International Conference on Artificial Intelligence and Statistics (AISTATS...; Conference. Peer-reviewed conference (PDF footer; Semantic Scholar) \\
\textbf{S24} \cite{S24} Lofstrom et al. (2024); RQ1, RQ3 & Local feature-importance explanations are unstable and distorted by poorly calibrated models, and feature-importance uncertainty is rarely reported. & Calibrated Explanations (CE): Venn-Abers calibration plus uncertainty intervals for each feature weight; factual and counterfactual rule-based feature weights; uses LIME's discretiser; random forest and XGBoost. & 25 public binary classification datasets from UCI and PROMISE (p.12). & Variance of feature weights over 30 explanations with only the seed changed (stability) and over 30 resampled refits (robustness); 20 stratified test instances. & CE was stable by construction (minute variation from seeds); robustness variance was low but clearly larger than stability variance (p.14). & Expert Systems with Applications; Journal. Peer-reviewed journal (Crossref) \\
\textbf{S25} \cite{S25} Pala \& García-Cuesta (2024); RQ1, RQ3 & Random-forest SHAP explanations can change sharply for tiny input changes because a single weak tree switches branch. & AXOM: discriminative averaging of weak-learner (tree-level) SHAP explanations to build a more robust ensemble explanation; robustness measured with a local Lipschitz estimate over neighbourhoods where \textbar{}x\_\allowbreak{}i,a - x\_\allowbreak{}j,a\textbar{} \textless{}= eps and f(x\_\allowbreak{}i) = f(x\_\allowbreak{}j) (eps = 0.01). & UCI Wine (13 features), Glass (10), Seeds (7), Banknote Authentication (5) (Table 1, p.5). & Local Lipschitz robustness (mean, SD) over 10,000 perturbed points per instance; baselines: decision-tree SHAP and standard random-forest SHAP. & Combining weak-learner explanations improved robustness relative to standard RF SHAP on most datasets; one exception where a decision tree appeared perfectly robust despite lower accuracy (pp.5-6). Shows that crossing a split in one tree makes the Lipschitz ratio diverge. & arXiv:2402.19025; Preprint. Preprint (not peer reviewed) \\
\textbf{S26} \cite{S26} Ribeiro et al. (2024); RQ1, RQ2 & How reliable and stable are XAI methods when test data are perturbed? & Pipeline with LGBM, MLP, decision tree and KNN; test data perturbed at increasing levels; feature-relevance ranks from Dalex, Eli5, eXirt, Lofo and SHAP compared; IRT-based eXirt reliability analysis. & Diabetes dataset (binary classification) (p.7). & Spearman correlation and bump charts of feature ranks across perturbation levels; Friedman-Nemenyi tests for model performance (p.9). & Most XAI methods, including SHAP, changed rankings under perturbation; eXirt identified less reliable models (Sec.5). & arXiv:2407.03108; Preprint. Preprint (not peer reviewed) \\
\textbf{S27} \cite{S27} Simpson et al. (2024); RQ2, RQ3 & Which explainer is most robust, and how does network smoothness relate to explanation robustness? & Lower bounds on probabilistic Lipschitzness of Integrated Gradients, LIME and SmoothGrad; normalised astuteness metric; link between a network's local Lipschitz constant and its stable rank. & XOR, Iris (tabular) and MNIST; MLPs with 2-16 layers (p.12). & Normalised astuteness of explainers; stable rank as heuristic. & Integrated gradients were less robust than SHAP or LIME; stable rank agreed with normalised astuteness (Sec.7). & arXiv:2402.18863; Preprint. Preprint (not peer reviewed) \\
\textbf{S28} \cite{S28} Bacevicius et al. (2025); RQ1, RQ3 & LIME's normal-distribution perturbations fit poorly to skewed, imbalanced network-traffic features, giving unstable explanations. & Modified LIME sampling with Weibull, Gamma, Beta and Pareto distributions for DT, RF, K-NN and XGBoost classifiers; fidelity (R2) and stability score S compared with default LIME. & CIC-IDS-2018 network-flow features (CIC-IDS-2017 discussed). & Local surrogate R2 (fidelity) and stability score S across repeated explanations. & Pareto gave the highest fidelity and stability, e.g. K-NN R2 = 0.9971, S = 0.9907 and DT R2 = 0.9267, S = 0.9797 (results). & Machine Learning and Knowledge Extraction; Journal. Peer-reviewed journal (Crossref) \\
\textbf{S29} \cite{S29} Ballegeer et al. (2025); RQ1, RQ2 & Whether instance-dependent cost-sensitive (IDCS) loss functions change the stability of SHAP and LIME explanations in credit scoring. & Traditional vs IDCS variants of XGBoost, logistic regression, random forest and neural network (8 models). Training sets resampled to default rates of 1-30\% with 25 resampling iterations per rate; SHAP and LIME feature-importance rankings recomputed on the test set each iteration. & Four public credit datasets: GMSC (112,915 x 10; 6.74\% default), VUB (18,667 x 20; 16.95\%), HMEQ (5,960 x 12; 19.95\%), South German Credit (1,000 x 20; 30\%) (Table 2, p.5). & Performance: AUC, AP, Brier, savings, relAEC (new). Stability: coefficient of variation (attribution values) and Sequential Rank Agreement (rankings). Friedman test with Hommel-adjusted post-hoc tests for performance (pp.13-14). Baseline: cost-insensitive counterpart of each model. & IDCS classifiers improve cost-efficiency but yield significantly less stable SHAP and LIME explanations than traditional models; stability falls as class imbalance rises for all four model types, more strongly for IDCS models (abstract; Sec.6). & European Journal of Operational Research; Journal. Peer-reviewed journal (Crossref) \\
\textbf{S30} \cite{S30} Fan et al. (2025); RQ3 & Monte Carlo Shapley approximations are slow and their estimates vary between runs. & SIM-Shapley: stochastic iterative momentum with adaptive mini-batches, l2 regularisation, negative-sampling detection and bias correction; proofs of linear Q-convergence and variance bounds. & Tabular datasets incl. Airbnb, Bank, Bike sharing and (South German) Credit (appendix figures). & Consistency (Pearson correlation across repeated runs vs a benchmark), bias vs ground truth, running time; baselines incl. KernelSHAP-type estimators and SAGE. & Reduced computation time by up to 85\% while maintaining or improving attribution quality and run-to-run consistency (abstract; eFigs.2-3). & arXiv:2505.08198; Preprint. Preprint (not peer reviewed) \\
\textbf{S31} \cite{S31} Goldwasser \& Hooker (2025); RQ1, RQ3 & Monte-Carlo sampling makes the top-ranked features from SHAP/\allowbreak{}LIME unstable across replicates. & Retrospective verification of stable ranks via multiple testing, and two sampling algorithms (StableSHAP/\allowbreak{}SPRT-SHAP, adjusted S-LIME) that identify the top-K set or ranking with bounded family-wise error rate. & Tabular benchmarks incl. Adult Census Income; neural-network classifiers (Fig.1, Sec.6). & Family-wise error rate of rankings/\allowbreak{}sets over repeated runs; number of verified stable ranks; compared with Holm-adjusted confidence intervals and plain Shapley sampling/\allowbreak{}LIME. & Plain Shapley sampling and LIME produced unstable top-5 rankings on Adult; the proposed procedures kept FWER near the nominal alpha (e.g., 0.1) while using computation adaptively (Fig.1; Sec.6). & Proc. 41st Conference on Uncertainty in Artificial Intelligence (UAI 2025), PMLR 286; Conference. Peer-reviewed conference (per Semantic Scholar); arXiv v5 read \\
\textbf{S32} \cite{S32} Hermosilla et al. (2025); RQ1 & Forensic use of IDS explanations needs reliable, reproducible SHAP/\allowbreak{}LIME explanations. & XGBoost and TabNet IDS models; SHAP and LIME local and global explanations; fidelity, stability and Jaccard similarity metrics. & UNSW-NB15 network-flow dataset. & Fidelity score, explanation stability, Jaccard similarity of top features. & XGBoost (97.8\% validation accuracy) outperformed TabNet in explanation stability and global coherence; SHAP and LIME were complementary (abstract). & Applied Sciences; Journal. Peer-reviewed journal (Crossref) \\
\textbf{S33} \cite{S33} Knab et al. (2025); RQ1 & The many LIME adaptations proposed to fix fidelity, stability and applicability problems are hard to navigate. & Survey of works that explicitly address LIME limitations or improvements; concept matrix; taxonomy by LIME sub-steps (feature generation, sample generation, attribution, representation) and key issues; continuously updated online overview (pp.4-5). & Not applicable (review). & Not applicable (review). & Stability is one of the main problems that LIME variants target; most variants are evaluated on only a subset of criteria, so it is unclear which is superior (p.16). & Explainable Artificial Intelligence (xAI 2025), Part II, Communications in Computer and...; Conference. Peer-reviewed conference (Crossref title match) \\
\textbf{S34} \cite{S34} Lin \& Wang (2025); RQ1, RQ2 & Consistency of SHAP explanations across re-trained probability-of-default models, motivated by adverse-action requirements. & 100 XGBoost models with fixed hyperparameters and different random seeds; global SHAP values and ranks compared across models. & Default of Credit Card Clients (Taiwan): 30,000 accounts, 23 raw features (79 after one-hot encoding); 22.1\% default (pp.2-3). & KS statistic, confusion-matrix metrics for performance; Kendall's W for SHAP-rank concordance across 100 models. & KS varied narrowly (41.31-43.63\%, median 42.37\%), yet SHAP stability depended on importance level: Kendall's W = 0.93 for the top five features vs 0.34 for six mid-importance features (abstract, Sec.3). & Risks; Journal. Peer-reviewed journal (Crossref) \\
\textbf{S35} \cite{S35} Nakanishi (2025); RQ1, RQ3 & Most XAI methods give single-point explanations without reliability information, and these explanations change with small data or seed changes. & Bayesian extension of approximate inverse model explanations (AIME) with closed-form posterior and credible intervals; bootstrap stability analysis against SHAP and LIME; synthetic validation of interval coverage. & Synthetic data with known importance; Titanic survival; breast-cancer diagnosis (p.1). & Spearman rank correlation and normalised attribution standard deviation across bootstrap resamples; CI coverage. & Kept mean Spearman rank correlation (0.969), improved explanation stability by 7\% and reduced normalised attribution SD by 58\% compared with SHAP and LIME (p.2). & IEEE Access; Journal. Peer-reviewed journal (Crossref) \\
\textbf{S36} \cite{S36} Rai et al. (2025); RQ1, RQ2 & How class imbalance affects the consistency of LIME, SHAP and PDP explanations for clinical risk models. & XGBoost, random forest and MLP trained on several imbalance ratios; explanations compared with a balanced (undersampled) reference model. & CPRD UK primary-care records for lung-cancer prediction (anonymised; restricted licence). & Jaccard similarity of top-10 features and Rank Agreement vs balanced reference; Mann-Whitney U (performance) and Wilcoxon signed-rank tests on 100 bootstrap samples (pp.5,7). & Explanation consistency degraded as imbalance increased, with differences between LIME and SHAP and across models; balanced training gave more consistent explanations (results). & Frontiers in Artificial Intelligence; Journal. Peer-reviewed journal (Crossref) \\
\textbf{S37} \cite{S37} Wang \& Kee (2025); RQ1 & Shapley attributions fluctuate under data shifts, input perturbations and adversarial attacks. & Weighted fusion of Shapley values and Least Core attributions (linear programming) for XGBoost; robustness tests under several data perturbations. & 2024 Beijing rental listings from Beike/\allowbreak{}Lianjia (80/\allowbreak{}20 split). & Explanation consistency under perturbations; fairness of attribution. & The hybrid preserved Shapley global fairness while significantly improving explanation consistency under perturbations (abstract). & Buildings; Journal. Peer-reviewed journal (Crossref) \\
\textbf{S38} \cite{S38} Yan et al. (2025); RQ1 & Single landslide-susceptibility models give unstable SHAP interpretations. & Heterogeneous hybrid of LightGBM, XGBoost and random forest interpreted with SHAP; 10-fold cross-validation; coefficient of variation (CV) of SHAP importances; McNemar tests for accuracy. & 347 landslide and 347 non-landslide points in Xi'an with conditioning factors (slope, elevation, LS factor, etc.) from a national landslide inventory 1949-2011. & CV of SHAP importance across folds (lower = more robust); accuracy, F1, AUC; McNemar's chi-squared test. & Hybrid CV 0.175 vs 0.208 (LightGBM), 0.240 (XGBoost) and 0.207 (RF) with comparable accuracy (abstract). & ISPRS International Journal of Geo-Information; Journal. Peer-reviewed journal (Crossref) \\
\textbf{S39} \cite{S39} Acampora \& Vitiello (2026); RQ1, RQ3 & Whether LIME explanations of quantum classifiers are stable enough to be trusted in critical domains. & Quantum SVM and quantum neural network classifiers (Qiskit), ideal and noisy simulation; LIME tabular explainer; repeated LIME calls per instance; Visani et al.'s VSI and CSI indices. & Tabular classification datasets from critical domains (UCI repository). & Variables Stability Index and Coefficients Stability Index (0-100) across repeated applications. & LIME gave explanations for quantum SVMs and QNNs that were interpretable and stable in both ideal and noisy quantum settings (abstract). & Machine Learning; Journal. Peer-reviewed journal (Crossref) \\
\textbf{S40} \cite{S40} Alhawarat et al. (2026); RQ1, RQ2 & SHAP explanations in crash-severity models are rarely audited for faithfulness and stability. & Five classifiers with temporal splits; isotonic calibration; four-step quantitative SHAP audit (sanity, deletion/\allowbreak{}insertion faithfulness, perturbation stability via local-Lipschitz style test with 10 perturbed copies, sigma = 0.01 of feature SD plus one binary toggle); cross-jurisdiction transfer. & US-Accidents (7.0 million records, 2016-2023) and UK STATS19 (about 1,010,000 records, 2016-2022). & Low-stability fraction of instances under perturbation; deletion faithfulness; ECE; transfer metrics. & Explanation stability failed at realistic perturbation levels: 54.3\% of instances were low-stability at sigma = 0.05, driven by spatial sparsity in geohash cells (abstract). & Applied Sciences; Journal. Peer-reviewed journal (Crossref) \\
\textbf{S41} \cite{S41} Alkan (2026); RQ1 & Class-imbalance treatments, calibration and explanation stability are rarely examined together in IoT intrusion detection. & Five imbalance strategies with XGBoost and random forest at three label granularities; training pools redrawn with ten seeds against a fixed test sample; TreeSHAP rankings compared between baseline (S0) and SMOTE (S1) models. & CICIoT2023 flow records (binary, category and attack-type levels). & Macro-averaged metrics, ECE; Friedman and Wilcoxon-Holm tests; Spearman correlation of TreeSHAP rankings and top-5 overlap (p.8). & Class weighting was best for XGBoost (Friedman chi2 = 32.240, p = 1.71e-6); TreeSHAP rankings were stable under resampling (Spearman 0.994, significant) (pp.13,16). & Journal of Expert Systems and Sustainable Development; Journal. Peer-reviewed journal (Crossref) \\
\textbf{S42} \cite{S42} Bhuiyan et al. (2026); RQ1, RQ2 & Stable SHAP explanations may not reflect performance-based feature reliance under distribution shift. & Logistic regression, linear SVC and random forest; cross-variant evaluation; SHAP stability via Spearman correlation and top-5/\allowbreak{}top-10 overlap with the Base-model ranking (fixed model and retrained models); permutation importance with repeated permutations; bootstrap. & Bank Account Fraud dataset suite: Base plus five biased variants (group-size, prevalence, separability, temporal shift). & Spearman rank correlation, top-k overlap, rank gaps between SHAP and permutation importance; classification metrics under imbalance. & SHAP rankings stayed relatively stable across variants, but permutation importance showed weaker and more variable reliance on several SHAP-ranked features, so stable explanations did not guarantee reliable feature reliance (abstract). & FinTech; Journal. Peer-reviewed journal (Crossref) \\
\textbf{S43} \cite{S43} Caraker et al. (2026); RQ1, RQ3 & Retraining the same model with a different seed can invert which correlated feature SHAP ranks as most important. & Proves that no feature ranking can be faithful, stable and complete under collinearity; quantitative bounds for gradient boosting, Lasso, neural networks and random forests; DASH (ensemble-averaged SHAP) as resolution; Lean 4 formalisation (305 theorems). & Survey of 77 public datasets; experiments across 11 datasets and 3 GBDT implementations; financial case studies (p.6). & Flip rates of within-group rankings, expected Spearman ranking stability, bootstrap 95\% CIs over 10 trials (p.18); SNR-based flip-rate formula (R2 = 0.94 on Breast Cancer, p.35). & 68\% of 77 datasets show attribution instability; for collinear pairs the ranking is close to a coin flip; the attribution ratio diverges as 1/\allowbreak{}(1-rho\^{}2) for gradient boosting but converges for random forests; DASH yields stable rankings (pp.5-15). & arXiv:2605.21492 (also Zenodo, doi 10.5281/\allowbreak{}zenodo.19468379); Preprint. Preprint (not peer reviewed) \\
\textbf{S44} \cite{S44} Chi et al. (2026); RQ1 & Feature-group fusion models lack explicit group-level attribution and attribution-stability guarantees. & SWFF-R: convex fusion of a full-feature predictor and group experts weighted by TreeSHAP importances; Rademacher bound; fixed-coefficient Lipschitz stability proof; multi-seed repetition (20 seeds) and model-confidence-set tests (appendix). & Seven real-world tabular regression datasets and a synthetic 500K stress test. & RMSE; group-attribution stability; statistical comparison procedures. & SWFF-R preserved predictive performance and gave stable group attributions with a provable Lipschitz bound under local input perturbations (abstract; Sec.3.1.9). & Mathematics; Journal. Peer-reviewed journal (Crossref) \\
\textbf{S45} \cite{S45} Conti et al. (2026); RQ1, RQ3 & Attribution scores alone do not show how well an explainer aligns with a reference hypothesis or how variable it is. & Weight-of-Evidence hypothesis-testing framework for feature-importance methods with theory linking WoE to attribution variance; LIME and SHAP explaining a random forest, 50 runs on 50 instances with top-k (k=3). & Titanic (domain-knowledge hypotheses) and a synthetic classification dataset with known informative features (pp.10-11). & WoE vs ground truth and vs model-derived hypotheses; coefficient of variation of WoE across runs. & LIME gave variable WoE across runs, whereas SHAP was deterministic (23 instances woe=+inf, 27 woe=-inf on Titanic); WoE estimates were consistent (80\% of CVs between 0 and 0.3) (pp.10,14). & arXiv:2609.00090 (accepted at XKDD 2026 workshop, ECML-PKDD, per arXiv comment); Preprint. Accepted, not yet published (authors' statement); preprint read \\
\textbf{S46} \cite{S46} Conti et al. (2026); RQ1, RQ3 & Stochastic components make attribution rankings vary between runs; how far down a ranking can be trusted? & Distribution-based framework: attribution distributions over repeated runs; pairwise separability of adjacent features (continuous Jaccard distance) defines the largest reliable rank (k-stability); explainers compared by the proportion of instances where one is more stable. & Diabetes, Heart Disease, Mobile and Churn tabular benchmarks (p.7). & k-stability over 50 instances x 50 runs per explainer (LIME, SHAP, DiCE, random baseline); +/\allowbreak{}-2 sigma/\allowbreak{}sqrt(n) intervals; sensitivity to number of runs (pp.7-10). & SHAP was quasi-deterministic (perfect stability) on Churn and Diabetes; LIME rankings were reliable only for the first few positions; the metric was not very sensitive to the number of runs (pp.8-11). & arXiv:2608.02697 (accepted at EXPLAINS 2026 per arXiv comment); Preprint. Accepted, not yet published (authors' statement); preprint read \\
\textbf{S47} \cite{S47} Dubey et al. (2026); RQ3 & SHAP values for tree ensembles are reported as point estimates, ignoring their variability. & UbiQTree: decomposes SHAP-value uncertainty into aleatoric, epistemic and entanglement components using Dempster-Shafer theory and Dirichlet-process hypothesis sampling over tree ensembles (random forests). & Several tabular health datasets incl. SEER Breast Cancer and a hospital length-of-stay task (pp.17-22). & Variance components, 95\% credible intervals, coefficient of variation of SHAP; correlation between importance and SHAP SD. & Epistemic uncertainty in SHAP explanations is substantial and linked to data quality and model development; in some datasets the most important features are also the most unstable (pp.1, Fig.15). & Patterns; Journal. Peer-reviewed journal (Crossref) \\
\textbf{S48} \cite{S48} Ibebuchi (2026); RQ1, RQ3 & Highly ranked SHAP predictors are not necessarily reliable explanations across spatial, resampling, model and predictor-correlation uncertainty. & Station-held-out validation of tree models; SHAP importance; station-level bootstrap refits; rank stability (Spearman, top-10 overlap), directional consistency, ablation and multi-model comparison combined into an RPI; controlled grouped benchmark. & Daily discharge at 51 USGS stations in the Chesapeake Bay watershed, 1990-2024, with time-safe lagged climate-mode predictors (station-season rows) (p.5). & Within-model rank, Spearman correlation, top-10 overlap; bootstrap CIs (pp.9-11). & RPI recovered the predefined stable drivers more accurately than raw SHAP importance in the benchmark; September PDO was the most spatially transferable and explanation-stable predictor (pp.1,22). & Machine Learning: Science and Technology; Journal. Peer-reviewed journal (Crossref) \\
\textbf{S49} \cite{S49} Kaden et al. (2026); RQ1 & Near-optimal models (Rashomon set) can give different explanations; consistency of explanations across them is unexplored. & Shallow two-layer MLPs trained repeatedly; exact Shapley values (low-dimensional data) and weight-based importance profiles; cluster analysis of importance profiles across the Rashomon set. & Pima Indians Diabetes, Wheat Seed, and a third small tabular dataset (Table 1, p.5). & Spearman correlation between importance profiles, number of clusters, adjusted Rand index between cluster solutions. & Correlations between importance profiles and agreement of cluster solutions were not very high (e.g., ARI 0.26, 0.092, 0.07), showing qualitatively different explanations among equally accurate models (Table 1). & ESANN 2026 proceedings; Conference. Peer-reviewed conference proceedings (Crossref) \\
\textbf{S50} \cite{S50} Mahendra \& Khare (2026); RQ1, RQ2 & Whether explanations stay stable when models are compressed by post-training quantisation, given the assumption that preserved accuracy implies preserved explanations. & Single FP32 logistic regression, random forest and XGBoost models quantised to INT8 and INT4 without retraining (for XGBoost, quantisation applied at the explanation layer); SHAP explanations compared with the full-precision model. & IEEE-CIS Fraud Detection (Kaggle): 420 engineered features; stratified split 354,324 /\allowbreak{} 118,108 /\allowbreak{} 118,108 (p.4). & Predictive metrics (incl. ROC-AUC); global feature-rank correlation, mean attribution drift and instance-level cosine similarity of SHAP vectors. & Predictive performance and explanation stability were not aligned: nearly identical ROC-AUC was accompanied by noticeable changes in global rankings and local attribution vectors; aggressive (INT4) quantisation caused substantial attribution drift (abstract; p.10). & Cureus Journal of Computer Science; Journal. Peer-reviewed journal (Crossref) \\
\textbf{S51} \cite{S51} Maseno et al. (2026); RQ1, RQ2 & High IDS detection performance does not show whether explanations are reliable. & XGBoost, logistic regression, random forest; SHAP vs permutation importance agreement; sufficiency/\allowbreak{}comprehensiveness retraining tests; retraining-stability analysis over three seeds; spurious-feature injection; radar profile. & UNSW-NB15 and IoT-ToN. & Spearman agreement; top-k faithfulness tests (k = 3, 5, 10); retraining stability. & Logistic regression gave the most stable explanations but weaker detection; high performance did not imply trustworthy explanations (abstract). & Journal of Computer Virology and Hacking Techniques; Journal. Peer-reviewed journal (Crossref) \\
\textbf{S52} \cite{S52} Mikriukov et al. (2026); RQ2, RQ3 & Explanations are costly and their reliability is not guaranteed; can epistemic uncertainty flag where explanations will be unstable? & Five architectures (incl. logistic regression, gradient-boosted trees with an RF surrogate for uncertainty, MLP) and four XAI methods (incl. SHAP and LIME); explanations recomputed under input perturbations (noise and BIM adversarial, epsilon-bounded); stability correlated with epistemic uncertainty; routing and gating use cases. & Four UCI tabular classification datasets (Table 1, p.6; 70/\allowbreak{}15/\allowbreak{}15 split, z-scored); PlantVillage images as a generalisation check. & Spearman's rho and Kendall's tau between original and perturbed attribution rankings; faithfulness measures; correlation with epistemic uncertainty. & Strong negative correlation between epistemic uncertainty and explanation stability across datasets, models and XAI methods; uncertainty also separates faithful from unfaithful explanations (abstract). & arXiv:2603.29915; Preprint. Preprint (not peer reviewed) \\
\textbf{S53} \cite{S53} Najih et al. (2026); RQ1, RQ3 & SHAP interpretation stability may depend on data characteristics and the model used. & XGBoost and LightGBM fitted repeatedly with fixed hyperparameters; 64 simulated datasets varying minority-class proportion, feature correlation and number of features; SMOTE/\allowbreak{}ADASYN on empirical data; fixed seeds. & Simulation population of 600,000 observations; anonymised BPJS Kesehatan claims 2018-2023 (22 features), 70:30 stratified splits. & Sequential Rank Agreement of SHAP rankings across repeated runs. & Higher feature correlation and extreme class imbalance reduced ranking stability; imbalance handling improved performance and interpretation stability on BPJS data (abstract). & Jurnal Statistika dan Komputasi; Journal. Peer-reviewed journal (Crossref) \\
\textbf{S54} \cite{S54} Nnadi et al. (2026); RQ1, RQ2, RQ3 & LIME's operating limits on large, high-dimensional behavioural data are unverified. & Benchmark of hierarchical LIME (H-LIME) vs LIME with logistic regression, random forest and XGBoost; 10 x 10 repeated cross-validation; probability calibration; 10 independent perturbation runs per instance. & Kaggle online-learning adaptability dataset (1,205 instances; mostly categorical) and OULAD behavioural logs (p.5). & Jaccard similarity and Spearman correlation across perturbation runs; Kruskal-Wallis tests for subgroup attributions (p.11). & LIME was stable on the socio-economic dataset (rho about 0.87) but collapsed on the high-dimensional behavioural dataset (rho about 0.35); H-LIME mitigated the instability (abstract). & IEEE Access; Journal. Peer-reviewed journal (Crossref) \\
\textbf{S55} \cite{S55} Nurhayati et al. (2026); RQ2 & Class imbalance may change LIME-based interpretations in nutritional-status prediction. & Random forest trained on original imbalanced data and on SMOTE-balanced data; LIME explanations compared between scenarios. & Nutritional-status (child health) tabular dataset. & Interpretation stability score from feature changes and contributions between scenarios. & Interpretation stability score reached 0.80: fairly stable but not fully robust to changes in data distribution (abstract; results). & Indonesian Journal of Infomatics; Journal. Peer-reviewed journal (Crossref) \\
\textbf{S56} \cite{S56} Padi et al. (2026); RQ1 & Knowledge-tracing models are opaque; the study tests whether TreeSHAP explanations of a feature-based model are stable and faithful. & XGBoost with TreeSHAP on 13 engineered behavioural features vs four deep KT baselines; stability across cross-validation folds, random seeds and two TreeSHAP conditioning schemes; retraining-based faithfulness ablations. & ASSISTments 2009 (rebuilt to remove leakage) and ASSISTments 2012; users with \textgreater{}=20 interactions (pp.5-6). & AUC; Spearman correlation of global rankings across 15 fold-by-seed runs; ablation of top-ranked features; student-level bootstrap for uncertainty. & Rankings were consistent across folds, seeds and conditioning schemes (Spearman 0.989-1.000); removing top TreeSHAP features hurt AUC more than random removal (abstract). & arXiv:2609.28502; Preprint. Preprint (not peer reviewed) \\
\textbf{S57} \cite{S57} Piyavechvirat et al. (2026); RQ1, RQ2 & Clinical ML benchmarks rarely quantify how SHAP attributions deform under noise, shift and surrogate drift. & Seven synthetic clinical 'worlds' (noise, missingness, outliers, shift, surrogate drift, label noise); linear, tree, kernel, neural and ensemble models over 15,000+ configurations; integrity profile R, S, C, E combined into an RSCE index. & Synthetic clinical cohort; MIMIC-IV mortality used for external alignment (pp.6,19). & E score from global-importance rank correlation, local SHAP cosine similarity and normalised Wasserstein distance between reference and perturbed worlds; Spearman/\allowbreak{}Kendall with bootstrap CIs and permutation p-values (p.18). & Correlations between synthetic RSCE and real-world degradation were modest and uncertain because few models were evaluated (p.18). & IEEE Access; Journal. Peer-reviewed journal (Crossref) \\
\textbf{S58} \cite{S58} Qi et al. (2026); RQ1 & Attrition models are tuned for discrimination, yet nothing requires their SHAP driver ranking to survive retraining. & Bi-objective hyperparameter optimisation (cost-sensitive AUPRC and expected agreement of SHAP global rankings across bootstrap refits); U-statistic estimator; minimum-variance ensemble weights; multi- vs single-objective solvers. & Five public HR-attrition benchmarks (1.4-50.6\% positives), incl. IBM-HR. & Expected ranking agreement across refits; AUPRC; statistical tests between solvers. & Attribution variance is not a proxy for explanation stability (minimising it picked the least stable learner on every benchmark); no test separated multi- from single-objective solvers at a practitioner's budget (abstract). & Mathematics; Journal. Peer-reviewed journal (Crossref) \\
\textbf{S59} \cite{S59} Rajhans \& Khawarey (2026); RQ1, RQ2, RQ3 & Adversarial attacks may destabilise SHAP explanations of tree-ensemble security classifiers even when predictions appear robust. & Random forest and XGBoost attacked with five attacks (incl. black-box ZOO, Square Attack, HopSkipJump, plus PGD); TreeSHAP drift summarised in a new Explainability Stability Index (ESI) on the same [0,1] scale as a Robustness Index (RI). Extends an earlier MLP study. & Four tabular security datasets: phishing URLs, UNSW-NB15, NF-ToN-IoT, HIKARI-2021. & RI (prediction robustness) and ESI (attribution drift) per attack and model; PGD step-size ablation. & ZOO gives degenerate results on XGBoost (RI about 0.98) because of piecewise-constant prediction surfaces, while Square Attack shows real vulnerability (RI about 0.36). Even with near-perfect ZOO robustness, XGBoost ESI is about 0.06-0.16 vs 0.14-0.29 for RF: prediction robustness and explanation stability are distinct axes (abstract). & arXiv:2607.01679; Preprint. Preprint (not peer reviewed) \\
\textbf{S60} \cite{S60} Rajhans \& Khawarey (2026); RQ2 & Adversarial perturbations degrade both the accuracy and the interpretability of cybersecurity classifiers. & Neural classifiers attacked with L-infinity FGSM and PGD; Robustness Index (area under the accuracy-perturbation curve); gradient-based feature sensitivity; SHAP attribution drift defined as the mean absolute change in SHAP values between clean and perturbed inputs (Eq.6) on 256 test samples; FGSM adversarial-training ablation (eps = 0.05). & Phishing Websites (11,000 URLs, 30 features) and UNSW-NB15 (175,000 flows, 42 features), both from Kaggle (p.3). & Accuracy vs eps; Robustness Index; per-feature mean absolute SHAP drift; no statistical tests. & Features with high gradient sensitivity also showed the largest SHAP drift, i.e. robustness and interpretability degrade together; adversarial training raised the Robustness Index by up to 9\% while keeping clean accuracy (abstract; Sec.4). & Proceedings of the 18th International Conference on Agents and Artificial Intelligence; Conference. Peer-reviewed conference proceedings (Crossref) \\
\textbf{S61} \cite{S61} Rasel et al. (2026); RQ1, RQ2 & Whether explanation stability differs between healthcare-payment and financial-transaction fraud models. & Logistic regression, random forest (200 trees) and gradient boosting (300 estimators); global SHAP importance per independent train-test split; fixed seed. & Kaggle Healthcare Provider Fraud Detection (CMS-derived, provider level) and ULB Credit Card Fraud (Dal Pozzolo et al. 2015; PCA features) (p.4). & Spearman rank correlation of mean \textbar{}SHAP\textbar{} rankings between splits; standard classification metrics. & Healthcare fraud models gave stable, meaningful explanations; financial fraud models were accurate but less stable (abstract; Sec.9). & Journal of Computer Science and Technology Studies; Journal. Peer-reviewed journal (Crossref) \\
\textbf{S62} \cite{S62} Tolay (2026); RQ1, RQ2 & IDS studies ignore explanation cost and local stability when judging explainable detection. & Leakage-safe CICIoT2023 corpus; LR, DT, RF, XGBoost; TreeSHAP cost measured; two perturbation protocols: adversarial and local prediction-preserving (accepted only if predicted class unchanged and confidence change \textless{}= 0.05) (p.5); selective-explanation policies. & CICIoT2023 (39 features) after deduplication and label-collision removal; natural and balanced test distributions. & Top-feature membership, attribution-rank, directional consistency and magnitude drift; TreeSHAP runtime; cost-coverage trade-offs. & At 5,000 samples TreeSHAP took 700.759 s for RF vs 1.471 s for XGBoost; RF showed the strongest base-level stability; XGBoost kept rank and direction but showed more top-feature turnover and magnitude drift (abstract). & arXiv:2608.10349; Preprint. Preprint (not peer reviewed) \\
\textbf{S63} \cite{S63} Tsai et al. (2026); RQ1 & Whether interpretations of football-performance models remain stable when transferred from elite to university football. & Random forest (TreeExplainer) and MLP (DeepExplainer) trained on elite event data and applied to a university domain; SHAP and a Counterfactual Impact Score; five random seeds (0-4). & Top-five European league match event data (WhoScored/\allowbreak{}Opta; 8:1:1 match-level split) and NTHU university match data (small, private) (pp.5-6). & Seed-wise Spearman rank correlation of SHAP importance; cross-domain structural agreement; MAE/\allowbreak{}RMSE for performance. & Seed-wise SHAP ranking stability was high in elite data (RF 0.932 +/\allowbreak{}- 0.019; MLP 0.826 +/\allowbreak{}- 0.050) but low in the university domain (RF 0.420 +/\allowbreak{}- 0.018; MLP 0.323 +/\allowbreak{}- 0.178) (p.14). & arXiv:2605.10796; Preprint. Preprint (not peer reviewed) \\
\textbf{S64} \cite{S64} Tzionis et al. (2026); RQ1, RQ3 & LIME explanations are unstable because different neighbourhood samples give different explanations. & AKDE-LIME: density-aware kernel weighting of perturbed samples combined with proximity; compared with LIME, KernelSHAP, TreeSHAP and Anchor on five tree-based models; grid of noise levels (0.05-0.20) and 10-50 evaluated instances. & A real-world tabular dataset from the UCI repository (Data Availability). & Stability (rank-correlation-based score after Zhou et al.) and robustness (relative change in importance after Alvarez-Melis and Jaakkola); Jaccard distance for Anchor rules. & AKDE-LIME was significantly more stable and robust than LIME and reached stability comparable to TreeSHAP (e.g., decision tree 0.0011 vs 0.0158) (Results; Conclusions). & Applied Artificial Intelligence; Journal. Peer-reviewed journal (Crossref) \\
\textbf{S65} \cite{S65} Vaduva et al. (2026); RQ1, RQ2 & Existing stability measures treat all features equally and lack business meaning; effect of class imbalance (SMOTE) on explanation stability was unexamined. & Proposes CIES, a rank-weighted (harmonic) distance between original and perturbed SHAP attributions under a 'business noise' neighbourhood; four tree-based classifiers; raw vs SMOTE-balanced training; LIME re-run to test method-agnosticism. & Telco Customer Churn (7,043 x 19), German Credit (1,000 x 20), IBM HR Attrition (1,470 x 31) (p.8). & CIES vs uniform (unweighted) baseline using Wilcoxon signed-rank tests over 24 configurations; sensitivity over four noise levels; Lipschitz comparison. & Model complexity affects explanation credibility; SMOTE changes explanation stability as well as performance; CIES discriminates better than the uniform baseline (p\textless{}0.01 in all 24 configurations) (abstract). & arXiv:2603.05024; Preprint. Preprint (not peer reviewed) \\
\textbf{S66} \cite{S66} Vourganas \& Michala (2026); RQ1, RQ3 & Effect of correlated (collinear) features on the stability of post-hoc explanations in intrusion detection. & Theorem that multicollinearity inflates attribution variance; dataset-level VIF/\allowbreak{}correlation audit; linear, tree-based, kernel and neural models on full vs pruned feature sets; Explainability Fragility Score; mitigations CAA-Filtering (grouping attributions) and SHARP (training-time attribution-stability regulariser). & UNSW-NB15 network-intrusion benchmark. & Kendall's tau across bootstrapped explanations; Explainability Fragility Score; predictive metrics before/\allowbreak{}after pruning and regularisation. & Explanations and feature importances are non-identifiable under multicollinearity; pruning and SHARP improve explanation stability while predictive performance stays stable (abstract). & arXiv:2605.22529; Preprint. Preprint (not peer reviewed) \\
\textbf{S67} \cite{S67} Waguih (2026); RQ1, RQ2 & Clinical AI decision support is judged mainly on predictive metrics, not on explanation stability and related properties. & Decision tree, SVM and XGBoost; 30 independent runs; five dimensions: predictive reliability (ROC-AUC), explainability (SHAP-LIME agreement), explanation stability, clinical plausibility, imbalance robustness (G-mean). & METABRIC HER2 status (1,980 cases: 1,733 negative, 247 positive). & Explanation-stability score across runs alongside the other dimensions. & XGBoost had the highest ROC-AUC (0.983) and explanation stability (0.652 +/\allowbreak{}- 0.045) (abstract). & Future Business Journal; Journal. Peer-reviewed journal (Crossref) \\
\textbf{S68} \cite{S68} Xiang et al. (2026); RQ1, RQ3 & SHAP values and rankings vary across train-test splits, seeds and fitting procedures but are reported as fixed. & Bootstrap resampling and refitting; kernel density estimation of SHAP-score distributions (with an asymptotic Gaussian approximation); RoSHAP combines activity, strength and stability into a robust ranking criterion. & Simulations and real datasets incl. high-dimensional gene-expression data (Fig.1) and an image example (Sec.5). & Ranking stability across repeated analyses; recovery of signal vs chance-selected features in simulations. & RoSHAP improves interpretation stability and separates consistent signals from chance-selected features (Sec.6). & arXiv:2605.15154; Preprint. Preprint (not peer reviewed) \\
\textbf{S69} \cite{S69} Zhao (2026); RQ1, RQ2 & How to evaluate explainable credit underwriting jointly on profit, fairness, recourse and explanation stability. & Logistic regression, random forest (400 trees) and gradient boosting; bad-rate-constrained threshold policies; SHAP explanations; counterfactual recourse; bootstrap (B=5) stability analysis. & FICO HELOC (10,459 applicants; 23 bureau features); 70/\allowbreak{}30 stratified split (p.2). & Approval rate, profit, DP/\allowbreak{}EO fairness gaps, recourse success; stability = Spearman rho of global SHAP and cosine similarity of local SHAP for 50 fixed test instances (Table 3). & Gradient-boosting SHAP explanations were stable (mean Spearman 0.871 global; mean cosine 0.856 local) (abstract). & Artificial Intelligence and Machine Learning Review; Journal. Journal article (Crossref); peer-review process not stated in the PDF \\
\textbf{S70} \cite{S70} Fan et al. (2021); RQ1 & Analysts cannot tell whether explanation approaches for malware classifiers can be trusted. & Three quantitative metrics (stability across models/\allowbreak{}runs, robustness to adversarial-sample changes, effectiveness) applied to five explainers on classifiers for malware detection and family identification; top-k feature overlap for k = 1-20. & Five widely used Android-malware datasets built in different years (binary feature vectors, e.g., permissions and API calls) (p.1). & Stability, robustness and effectiveness metrics; comparison across explainers and learning algorithms. & Stability ranking LIME \textgreater{}= SHAP \textgreater{} Anchor \textgreater{} LORE \textgreater{} LEMNA; LIME and SHAP stability about 0.85; stability depends mainly on the explainer rather than the classifier (p.9). & IEEE Transactions on Information Forensics and Security; Journal. Peer-reviewed journal (Crossref) \\
\textbf{S71} \cite{S71} Ning et al. (2022); RQ3 & Explaining only the final model ignores the variability of SHAP importance across equally good models. & Generate nearly optimal models (loss-based Rashomon set) for an MLP by fine-tuning on diversified training subsets; SHAP global importance for each; uncertainty intervals and overall importance across models. & Electronic health records of 46,318 emergency-department patients admitted to Singapore General Hospital in 2017 (p.7). & Distribution and uncertainty intervals of mean \textbar{}SHAP\textbar{} across nearly optimal models. & Variable importance varied across nearly optimal models, and ShapleyVIC intervals gave a more robust picture than the single final model (abstract; Sec.3). & arXiv:2212.08370; Preprint. Preprint (not peer reviewed) \\
\textbf{S72} \cite{S72} Ranjbar \& Safabakhsh (2022); RQ1, RQ3 & LIME-type explanations of neural networks need better stability, fidelity and interpretability. & tree-ALIME: autoencoder-based LIME (ALIME) with a decision tree instead of a linear local model; 10,000 Gaussian perturbations per instance; stability score on top positive/\allowbreak{}negative coefficients over repeated runs; human interpretability test. & Four tabular datasets (incl. a liver-patients dataset) and a sentiment text task (pp.2-6). & Stability score, local fidelity, human prediction accuracy from explanations; comparison with ALIME. & tree-ALIME gave significantly better stability and local fidelity than ALIME and improved interpretability (abstract; Sec.IV). & 2022 27th International Computer Conference, Computer Society of Iran (CSICC); Conference. Peer-reviewed conference proceedings (Crossref) \\
\textbf{S73} \cite{S73} Cifuentes et al. (2024); RQ3 & SHAP scores depend on an assumed feature distribution that is usually estimated, so rankings may be misleading. & Distributional uncertainty modelled as product distributions in a hyperrectangle; SHAP-score intervals and complexity results; case study with sampling at different proportions (5 repetitions). & California Housing (20,640 samples, 8 features, binarised) with a depth-limited decision tree (pp.5-6). & SHAP intervals per feature; whether rankings change across the uncertainty region. & Feature rankings stayed informative when the distribution uncertainty was small but SHAP intervals widened with fewer samples (p.6). & ECAI 2024 (27th European Conference on Artificial Intelligence), Frontiers in Artificial...; Conference. Peer-reviewed conference proceedings (Crossref) \\
\textbf{S74} \cite{S74} Ketenci et al. (2024); RQ3 & Point-estimate Shapley attributions ignore data noise, predictive uncertainty and their effect on attributions. & PSI: variational masking-based network that jointly learns a predictor and latent stochastic Shapley values (mean = Shapley value) with credible intervals. & Synthetic data and 8 real-world datasets (e.g., ICU, census) (p.15); MNIST encoded as patches (p.12). & Predictive RMSE/\allowbreak{}PR-AUC vs baselines; recovery of true attributions and calibration of credible intervals in synthetic data. & PSI matched strong baselines in prediction while learning attribution distributions whose intervals covered the true attributions in synthetic tests (pp.11,15). & arXiv:2402.04211; Preprint. Preprint (not peer reviewed) \\
\textbf{S75} \cite{S75} B (2026); RQ2, RQ3 & Practice assumes that models with equivalent predictions give equivalent explanations. & Large empirical study of prediction-equivalent models from different hypothesis classes (tree-based, linear, neural); Spearman disagreement of SHAP (and LIME) vectors; Agreement Gap theorem linking disagreement to interaction density and axis-aligned tree splits (TreeSHAP interaction lemma); reliability score R(x). & 24 tabular datasets (MNIST only as a preliminary appendix extension) (pp.2,7). & Lottery rate (instances with substantial disagreement, rho \textless{} 0.5), within- vs cross-class Spearman, Cohen's d; correlation with dimensionality. & Prediction-equivalent models gave substantially different attributions; within-class rho = 1.000 vs cross-class 0.369 under split invariance; tree-linear agreement fell to 0.287 for d \textgreater{} 50 (r = -0.251, p \textless{} 0.001) (p.7). & arXiv:2603.15821; Preprint. Preprint (not peer reviewed) \\
\textbf{S76} \cite{S76} Bouyahia \& Kalogeratos (2026); RQ1, RQ3 & Simple LIME surrogates are stable but cannot follow complex non-linear decision boundaries. & NDT-LIME: neural decision trees as LIME surrogates; stability = mean pairwise cosine similarity of explanations over R random perturbation runs; regularity across nearby real data points. & Several benchmark tabular datasets (pp.1,12). & Stability, fidelity and regularity scores vs linear and decision-tree surrogates. & NDT-LIME balanced stability and fidelity, improving fidelity while its expressiveness introduced small variations between close neighbours (Sec.4; Conclusion). & arXiv:2603.20919; Preprint. Preprint (not peer reviewed) \\
\textbf{S77} \cite{S77} Breeden (2026); RQ2, RQ3 & LIME/\allowbreak{}SHAP-type explanations are unstable near decision boundaries, which is criticised in regulated credit settings. & Out-of-sample (OOS) Lipschitz and Hessian instability measures for local linear explanations and OOS Shapley values; conformal forecast-uncertainty; simulations with three data-generating functions and three model types (incl. gradient boosted trees, piecewise linear regression, sigmoid neural networks). & Simulated data (three generators). & Instability conditional on forecast-uncertainty region; one-sided Mann-Whitney tests (p.13). & Instability was stochastically lower in low-uncertainty regions (p near zero in all cases); piecewise-linear models such as gradient boosted trees showed the same instability near segment boundaries as non-linear models (pp.2,20). & arXiv:2602.00179; Preprint. Preprint (not peer reviewed) \\
\textbf{S78} \cite{S78} Sidhu (2026); RQ1 & Which post-hoc explanation method is most faithful, stable and comprehensible for ICU mortality prediction. & Gradient-boosting mortality models; TreeSHAP, ensemble LIME and Anchors compared on faithfulness (AOPC), stability (Lipschitz continuity estimate) and comprehensibility (explanation entropy/\allowbreak{}complexity). & MIMIC-IV ICU admissions (restricted access). & AOPC, Lipschitz estimate, explanation complexity, runtime. & TreeSHAP had the best faithfulness (AOPC 0.276) and stability (Lipschitz estimate 0.312), with 4 ms per explanation; Anchors were most comprehensible (abstract). & International Journal of Computer Sciences and Engineering; Journal. Peer-reviewed journal (Crossref) \\
\textbf{S79} \cite{S79} Xia et al. (2026); RQ1, RQ2 & Whether finer input data that improve predictive accuracy also improve explanation reliability. & Ten DEM-derived conditioning factors at 1-90 m resolutions aggregated to slope units; LR, SVM, RF, XGBoost under repeated random seeds; SHAP explanations scored for faithfulness, stability (attribution distance under Gaussian perturbation) and geomorphic consistency. & Event-based landslide inventory with slope units and DEMs resampled from a 1 m GF-7 DEM (pp.4-5). & AUC-type performance; faithfulness, stability and geomorphic consistency across resolutions and models. & Predictive performance and explanation reliability did not improve together; stability responded to resolution and model, and coarse DEMs could look stable by smoothing terrain while shifting feature-importance structure (abstract; p.7). & IEEE Journal of Selected Topics in Applied Earth Observations and Remote Sensing; Journal. Peer-reviewed journal (Crossref) \\
\end{longtable}
